\chapter{Literature Review}
\label{ch:literature-review}

\section{Cultural Appropriateness in Large Language Models}
Large language models are first trained to model natural-language distributions and are then commonly post-trained to better follow instructions and satisfy behavioral preferences. Reinforcement learning from human feedback (RLHF) is a prominent example: demonstrations and ranked outputs are used to train a preference signal that later guides policy optimization \parencite{ouyang2022instructgpt}. The important point for this thesis is that alignment is mediated by operational signals. Demonstrations, constitutions, human preferences, policies or evaluators define what behavior is rewarded.

Such signals are necessarily incomplete. A model can be generally helpful while still producing advice that is inappropriate for the explicit cultural, linguistic, legal or interpersonal context of a request. Conversely, disagreement with a population majority is not automatically culturally inappropriate. Pluralistic-alignment research therefore argues that systems serving heterogeneous populations may need to represent or accommodate multiple reasonable perspectives rather than collapse all judgments into one global preference \parencite{sorensen2024pluralistic}.

This thesis uses cultural alignment in an operational sense: a response should fit the explicit situation, avoid materially misrepresenting relevant practices or constraints, and preserve legitimate variation when the context does not justify a single categorical rule.

\label{sec:cultural-appropriateness}

The term \emph{cultural correctness} suggests a fixed, homogeneous cultural truth. This is too strong for many real-world questions. UNESCO describes culture through material, intellectual and emotional features, lifestyles, value systems, traditions and beliefs, while the UN Committee on Economic, Social and Cultural Rights treats cultural life as broad, evolving and connected to identity, language, knowledge, customs, institutions and participation \parencite{unesco2001culturaldiversity,cescr2009gc21}. Surveys of cultural NLP likewise note that research operationalizes selected proxies and facets rather than one universally accepted computational definition \parencite{adilazuarda2024culture,liu2025culturallyaware,pawar2025culturalawareness}.

Accordingly, this thesis uses \emph{cultural appropriateness} for a context-sensitive evaluative construct. Several epistemic distinctions matter:

\paragraph{Descriptive tendencies.} Evidence that a practice is common does not establish a universal rule. Cross-national opinion data can show population differences while preserving substantial within-population variation \parencite{durmus2023globalopinionqa,tao2024culturalbias}.

\paragraph{Social norms and etiquette.} Social expectations vary with relationship, setting, generation and locality. NormAd demonstrates how model judgments change as cultural norms become more explicit \parencite{rao2025normad}.

\paragraph{Pragmatic conventions.} Register, forms of address, indirectness, idioms and implicature affect cultural fit even when grammar and literal factuality are correct.

\paragraph{Values.} Values are plural and contested. Cultural evaluation should not transform majority preferences into moral unanimity.

\paragraph{Law and institutional rules.} Formal obligations differ from informal practice. SafeWorld demonstrates the need to consider both culturally situated norms and geographically varying legal constraints \parencite{yin2024safeworld}.

These distinctions motivate a verifier that matches the strength and kind of its conclusion to the available context and evidence.

\section{Cultural Evaluation and Benchmarking}
Recent surveys emphasize that culture has no single established computational operationalization. \textcite{adilazuarda2024culture} show that cultural NLP studies use heterogeneous proxies such as nationality, language, values, customs, food, traditions and social practice. \textcite{liu2025culturallyaware} similarly distinguish culturally aware NLP from culturally adapted NLP and propose a taxonomy covering cultural knowledge, language, values, norms and adaptation strategies. \textcite{pawar2025culturalawareness} further broaden the scope beyond text and emphasize the distinction between multilingual capability and cultural awareness.

This diversity is methodologically important. A factual multiple-choice benchmark about local customs, a value survey comparing population distributions and a culturally situated safety benchmark do not measure the same construct. Cultural evaluation should therefore make explicit what aspect of culture is under examination, which context is assumed, what evidence supports the judgment, and whether disagreement represents model error or legitimate cultural pluralism.

Population-level evaluation illustrates the problem. GlobalOpinionQA compares model outputs with cross-national opinion distributions rather than assuming one universal correct value response \parencite{durmus2023globalopinionqa}. \textcite{tao2024culturalbias} likewise report systematic cultural biases and show that cultural prompting can shift model outputs. Such steering is not automatically evidence of better cultural alignment, because prompting can also induce stereotyped or overly homogeneous portrayals of a requested population. Vericult therefore treats variation and uncertainty as part of the judgment rather than rewarding the mere presence of culturally specific wording.

BLEnD evaluates everyday cultural knowledge using question--answer pairs covering multiple countries, regions and languages \parencite{myung2024blend}. Its focus on daily practices is important because many culturally consequential failures concern food, routines, etiquette and local conventions rather than encyclopedic facts. CANDLE similarly extracts cultural commonsense knowledge at scale and organizes knowledge around contextualized practices and cultural facets \parencite{nguyen2023candle}.

CulturalBench moves toward a human-written and human-verified benchmark. The ACL 2025 dataset contains 1,696 questions across 45 regions and 17 topics, produced through a human--AI red-teaming process with multiple annotators \parencite{chiu2025culturalbench}. Its contribution is especially relevant because it demonstrates the difficulty of obtaining culturally meaningful evaluation material from model generation alone. Cultural benchmarks benefit from explicit human verification and from preserving cases in which more than one answer can be defensible.

These benchmarks establish that LLMs can be tested systematically for culturally situated knowledge, but they remain benchmark-specific. They do not directly provide a mechanism for evaluating an arbitrary free-text response after it has already been generated.

NormAd focuses on social norms rather than cultural trivia. It evaluates model behavior across situations in which the relevant norm is available at different levels of specificity, from broad value information to explicit cultural norms \parencite{rao2025normad}. The work is important for Vericult because it demonstrates that abstract cultural associations are not sufficient for reliable situational judgment. The same broad value can support different behavior depending on relationship, setting and local convention.

SafeWorld extends situated evaluation into safety and law. It contains geo-diverse queries spanning cultural norms and legal or policy constraints across many countries and regional settings \parencite{yin2024safeworld}. It highlights a distinction central to the D01--D10 framework: an informal norm, a public policy and a binding legal requirement are different epistemic objects. A verifier should not infer legal obligation from cultural prevalence or reduce a legal constraint to etiquette.

CultureBank takes a different approach by collecting community-derived cultural knowledge and contextualized cultural descriptions \parencite{shi2024culturebank}. This provides evidence that lived cultural knowledge cannot always be represented faithfully through official or encyclopedic sources. It also motivates Vericult's use of source type as provenance metadata rather than as a single fixed authority ranking.

FrameNet-Cultures evaluates open-ended generation through cross-cultural frame semantics. It defines cultural frames across multiple countries and studies outputs from contemporary model families followed by manual annotation \parencite{jamshidi2026framenet}. This demonstrates that culturally situated evaluation should extend beyond closed-form questionnaires.

ExCAM is even closer to Vericult's intended interface. It introduces an explainable cultural-awareness metric for arbitrary instruction--output pairs and constructs ExCAM40k from multiple cultural benchmarks with controlled cultural errors \parencite{leiter2026excam}. ExCAM identifies, rates and explains cultural errors rather than merely asking for a relative preference.

The architectural distinction is important. ExCAM is a learned metric; Vericult is an explicit multi-stage verification framework. Vericult's contribution is therefore not the discovery that free-text cultural errors can be classified. It is the investigation of whether applicability and assessability gates, dimension planning, deterministic target grounding, evidence retrieval and freezing, selective fallback/abstention and full traces provide useful additional structure.

\section{Cultural Alignment and Adaptation}
CultureLLM adapts language models to selected cultures using World Values Survey material, semantic augmentation and fine-tuning \parencite{li2024culturellm}. Its objective is to make the generator itself more culturally aligned. ValuesRAG instead retrieves culturally and demographically relevant values during generation and conditions the answer on this material \parencite{seo2025valuesrag}. Both demonstrate that external cultural information can improve model behavior.

Vericult differs in the direction of information flow. It does not retrieve cultural information to help the model write the answer. It receives an already generated response and asks whether that response is appropriate. This distinction matters for independence: if the same retrieved information is used to generate and evaluate a response, evaluation can become circular. Vericult therefore keeps its retrieval and evidence trace separate from the generator under evaluation.

\section{Model-Based Evaluation}

\subsection{Reward Models and Cultural Reward Modeling}
Reward models transform preference data into a reusable evaluation signal. A classifier RM often assigns a scalar reward to a prompt--response pair, while generative evaluators produce scores or judgments through language generation. General RM benchmarks such as RewardBench measure preference discrimination rather than cultural evidence gathering \parencite{lambert2025rewardbench}. Skywork Reward V2, used later as an independent baseline, is a general preference-model family developed through large-scale preference-data curation \parencite{liu2026skyworkrewardv2}.

A scalar reward is useful for ranking but limited diagnostically. It does not reveal which cultural claim, pragmatic convention, institutional constraint or stereotype caused the ranking. It is also a learned preference signal rather than a calibrated probability of cultural correctness. For this reason, the reward model remains strictly external to Vericult: RM scores never enter Vericult retrieval, scoring or ranking.

\label{sec:carb}

The closest prior work to the Best-of-4 component of this thesis is CARB, the Cultural Awareness Reward-model Benchmark \parencite{zhang2026carb}. CARB argues that general reward-model benchmarks do not adequately test whether reward models recognize culturally appropriate responses. The benchmark therefore constructs Best-of-N sets in which one response should be preferred over culturally problematic alternatives.

CARB covers ten culture/language settings and four broad domains: cultural commonsense, values, safety and linguistics. Its Best-of-N formulation is particularly useful for the present thesis because it supplies a natural comparative question for Vericult: can a post-hoc cultural verifier identify the culturally preferable answer among alternatives?

CARB also evaluates different reward-model families and reports that generative reward models can outperform conventional scalar classifier reward models on cultural preference. Its Think-as-Locals approach improves a generative reward model by eliciting culture-specific reasoning before scoring \parencite{zhang2026carb}. This supports a central hypothesis behind Vericult: culturally situated evaluation may benefit from explicit reasoning structure rather than a single opaque preference score.

However, CARB and Vericult solve different problems. CARB is fundamentally a reward-model evaluation and improvement study. It asks which candidate should be preferred in a Best-of-N setting. Vericult is designed to evaluate one arbitrary answer and expose a trace of why the answer is judged culturally appropriate or inappropriate. Best-of-4 is therefore retained as a comparison mode, not as the primary application.

A second difference is evidence. Think-as-Locals relies on the reward model's internal cultural knowledge and reasoning. Vericult additionally attempts to ground externally checkable claims in retrieved evidence and records whether that evidence was sufficient, conflicting or unavailable. The frozen architecture distinguishes a non-cultural task, a non-assessable answer and an evidence-insufficient cultural case. For unresolved context-dependent recommendations only, a narrow contextual fallback may judge the observable cultural handling of the advice; unresolved external facts and descriptive norms remain eligible for abstention.

For the final empirical study, CARB therefore has two roles. First, it motivates the reward-model comparison and Best-of-N setting. Second, where official instance-level CARB data are available, a deterministic subset is used for external validation after Vericult is frozen. Published CARB results are contextual comparison, not training labels or inputs to Vericult.

RewardBench provides a general framework for evaluating whether reward models discriminate preferred from rejected responses \parencite{lambert2025rewardbench}. The Skywork Reward V2 family builds on the broader reward-modeling literature and emphasizes large-scale preference-data curation \parencite{liu2026skyworkrewardv2}. These systems are valuable baselines precisely because they are not specialized implementations of Vericult.

In the thesis experiment, the reward model receives the same prompt--candidate pairs but remains independent of the Vericult pipeline. Its scores are converted into within-prompt rankings for Best-of-4 comparison. No reward score enters Vericult's dimension planner, retrieval, evidence memo or final score. This separation permits a meaningful comparison between a general learned preference signal and a structured cultural-verification process.

\subsection{LLM-as-a-Judge}
A second evaluation paradigm prompts an instruction-tuned LLM directly to compare or rate responses. LLM-as-a-judge evaluation is attractive because it is flexible and can express criteria in natural language. However, strong judges still exhibit systematic effects. Work on MT-Bench and related settings reports position, verbosity and self-enhancement biases, while later analyses show that response ordering can materially change pairwise judgments \parencite{zheng2023llmjudge,wang2024notfairevaluators}.

This thesis therefore treats a direct judge as a comparator rather than as the verifier. The direct judge sees the candidate responses and decides among them. Vericult instead decomposes the response, retrieves candidate-blind evidence where appropriate, and preserves the evidence path. Whether this additional structure helps is an empirical question rather than an assumption.

Direct LLM judging offers a second strong baseline. \textcite{zheng2023llmjudge} demonstrate that capable LLMs can serve as useful evaluators for general assistant responses, while also documenting position, verbosity and self-enhancement biases. \textcite{wang2024notfairevaluators} provide further evidence that response order can affect comparative judgments.

The direct judge in this thesis therefore uses deterministic candidate permutation and the same high-level D01--D10 rubric but no retrieval. This provides a stringent comparison: if a single judge call performs as well as the full verifier, the additional decomposition and evidence machinery must justify itself through robustness, diagnostics or external validity rather than through complexity alone.

The direct-judge literature also prevents an overly strong claim about Vericult. The verifier still uses an LLM for semantic decisions and remains vulnerable to model priors and prompting effects. Candidate-blind evidence synthesis and deterministic identifiers reduce particular failure modes; they do not eliminate evaluator bias.

\section{Evidence-Grounded Evaluation and Verification}
Retrieval-augmented generation established the broader principle that a language model can combine parametric knowledge with explicitly retrieved external information \parencite{lewis2020rag}. Vericult does not use retrieval to generate the user's answer; it uses retrieval to make parts of the evaluation inspectable. This is particularly valuable for claims about law, institutions, current policy, regional practice or the prevalence of cultural conventions.

Retrieval is not a source of ground truth. Search availability is uneven across languages and communities; provider snippets can omit qualifications; and source authority depends on the question. An official source is appropriate for a legal rule, while community material can be informative for lived practice. Vericult therefore uses source classes as provenance descriptors rather than a universal numeric hierarchy.

Retrieval can also contaminate evaluation if benchmark answers, annotations or result files are encountered. Benchmark-contamination research motivates explicit exclusions and artifact recording \parencite{deng2024contamination}. The implementation filters known repository and annotation paths and archives the exact evidence used, while acknowledging that copied or mirrored material cannot be eliminated perfectly.

\paragraph{Human reference versus ground truth.} Cultural annotations provide an operational comparison target, not objective universal truth. Agreement among annotators is itself part of the result.

\paragraph{Absolute and relative evaluation.} A candidate can be the best among four poor responses. Candidate-level appropriateness and Best-of-4 preference therefore answer different questions and must be analyzed separately.

\paragraph{Selective outcomes.} Vericult distinguishes several reasons not to return a substantive cultural label. A prompt can be \texttt{not\_culturally\_applicable}; an applicable prompt can receive a \texttt{not\_assessable} response; an assessable response can remain \texttt{insufficient\_evidence}; and individual dimensions may \texttt{abstain}. These states are reported explicitly rather than collapsed into one generic uncertainty value.

\paragraph{Paired comparison.} Because the same items are evaluated by Vericult, the reward model and the direct judge, system comparisons are paired at prompt level.

\paragraph{Reproducibility.} Fixing temperature is insufficient. Model versions, prompt templates, retrieval dates, evidence, ordering and code revisions can change outcomes. Vericult therefore distinguishes LIVE evidence acquisition from REPLAY evaluation over frozen retrieval artifacts.

These foundations motivate the research gap addressed in the next chapter: existing evaluators provide valuable preference and capability measurements, but a practical cultural verifier should additionally expose what is being judged, what external support was found and where its decision came from.

\section{Research Gap}
The literature provides strong components but distributes them across different objectives. CulturalBench, BLEnD and CANDLE evaluate cultural knowledge. NormAd and population-based work study contextual norms and value variation. SafeWorld integrates culturally situated safety with legal constraints. CultureLLM and ValuesRAG improve generation. CARB investigates cultural awareness in reward models. FrameNet-Cultures and ExCAM bring evaluation closer to open-ended free text.

The gap addressed here is deliberately narrower: an independent post-hoc verifier for an arbitrary prompt--response pair that combines:
\begin{enumerate}
    \item an explicit multidimensional cultural rubric;
    \item context extraction without unsupported demographic inference;
    \item an explicit cultural-applicability gate before D01--D10 planning;
    \item a response-assessability gate separating refusals/non-answers from cultural uncertainty;
    \item bounded selection of materially relevant response spans with Python-owned exact grounding and deterministic retrieval routing;
    \item retrieval only when the target is externally checkable;
    \item candidate-blind evidence synthesis after neutral verification questions have been formed;
    \item immutable evidence snapshots and traceable target--evidence links;
    \item explicit abstention plus a narrowly scoped recommendation fallback when external evidence remains unresolved; and
    \item both absolute single-response assessment and optional Best-of-4 ranking with an absolute acceptability gate.
\end{enumerate}

This is not a claim that each component is individually novel. The contribution is the combination and empirical evaluation of these properties as one standalone verifier.
